% !TEX program = xelatex
% Supporting Information for the 2026 Bundibugyo ebolavirus disease outbreak manuscript.
% Compiled standalone with the standard `article` class (PLOS accepts the SI as a
% separate PDF at submission).
\documentclass[11pt]{article}
\usepackage[margin=1in]{geometry}
\usepackage{booktabs}
\usepackage{url}
\usepackage[colorlinks=true,linkcolor=blue,citecolor=blue]{hyperref}
\usepackage{xcolor}

\title{Supporting Information for: Transmission dynamics of the 2026
Bundibugyo ebolavirus disease outbreak in the Democratic Republic of the
Congo: a compartmental model with an explicit observation layer}

\author{Tang Shujiang\\
School of Science, Hunan Institute of Science and Engineering, Yongzhou,
China. \texttt{sjtang@huse.edu.cn}}

\date{}

\begin{document}
\maketitle

\section*{S1. Verified weekly surveillance series (wk21--wk41, 2026)}
\label{sec:si_weekly}

The series below is the verified UKHSA \emph{Outbreaks under monitoring}
weekly counts at their DRC data dates, anchored at the 21 May 2026 WHO
Disease Outbreak News declaration. Values are confirmed cases and confirmed
deaths unless noted. Weekly increments are differences of consecutive data
dates (6--8 days apart). The three intervals 13 July--3 August (data dates
$t=67$, $74$) are merged into a single 21-day interval in the likelihood; the
$t=67\to74$ increment contains a 22 July retrospective reconciliation
(272 cases, 174 deaths), flagged with $\dagger$.

\begin{center}
\begin{tabular}{lrrrrr}
\toprule
Week ending & $t$ (d) & Cum.\ confirmed & Incr. & Cum.\ deaths & Death incr. \\
\midrule
2026-05-21 & 0   & 83   & --   & 9    & -- \\
2026-05-26 & 5   & 121  & 38   & 17   & 8  \\
2026-06-02 & 12  & 363  & 242  & 62   & 45 \\
2026-06-08 & 18  & 598  & 235  & 115  & 53 \\
2026-06-15 & 25  & 837  & 239  & 196  & 81 \\
2026-06-22 & 32  & 1,094& 257  & 277  & 81 \\
2026-06-29 & 39  & 1,333& 239  & 399  & 122\\
2026-07-06 & 46  & 1,708& 375  & 580  & 181\\
2026-07-13 & 53  & 2,011& 303  & 754  & 174\\
2026-07-20 & 60  & 2,473& 462  & 999  & 245\\
2026-07-27 & 67  & 3,360& 887$^\dagger$ & 1,487 & 488$^\dagger$ \\
2026-08-03 & 74  & 3,874& 514  & 1,751& 264\\
2026-08-10 & 81  & 4,449& 575  & 2,061& 310\\
2026-08-16 & 87  & 5,021& 572  & 2,378& 317\\
2026-08-24 & 95  & 5,656& 635  & 2,715& 337\\
2026-08-31 & 102 & 6,186& 530  & 3,007& 292\\
2026-09-07 & 109 & 6,757& 571  & 3,267& 260\\
2026-09-13 & 115 & 7,258& 501  & 3,510& 243\\
2026-09-20 & 122 & 7,773& 515  & 3,759& 249\\
2026-09-27 & 129 & 8,116& 343  & 3,924& 165\\
2026-10-04 & 136 & 8,603& 487  & 4,148& 224\\
2026-10-11 & 143 & 9,066& 463  & 4,381& 233\\
\bottomrule
\end{tabular}
\end{center}
\vspace{4pt}
\noindent $^\dagger$The $t=67\to74$ increment contains a 22 July 2026 retrospective
reconciliation (272 cases, 174 deaths).

\section*{S2. Data-source inventory with URLs}
\label{sec:si_sources}

All inputs are public surveillance series. The weekly case/death series
(S1) were transcribed from the UKHSA \emph{Outbreaks under monitoring}
reports and reconciled with WHO, ECDC, Africa CDC and DRC Ministry of
Health communications. The exact retrieval URL for each data date is
enumerated in the deposited repository and is provided at the
agency-portal level below.

\paragraph*{Surveillance and agency sources.}
\begin{itemize}
  \item UK Health Security Agency (UKHSA), \emph{Outbreaks under monitoring} weekly reports (wk21--wk41, 2026): \url{https://www.gov.uk/government/organisations/uk-health-security-agency}
  \item World Health Organization, Disease Outbreak News (declaration, 21 May 2026): \url{https://www.who.int/emergencies/disease-outbreak-news}
  \item European Centre for Disease Prevention and Control (ECDC), Ebola situation pages (hospitalized-in-isolation series): \url{https://www.ecdc.europa.eu/en/publications-data}
  \item Africa Centres for Disease Control and Prevention (Africa CDC), outbreak briefs: \url{https://africacdc.org/}
  \item DRC Ministry of Public Health (Minist\`ere de la Sant\'e), situation reports: \url{https://www.ministesante.gouv.cd/}
\end{itemize}

\paragraph*{Contextual and epidemiological sources cited in the manuscript.}
\begin{itemize}
  \item Villa J, Wilkinson A, Bedford J, Schmidt-Sane M, Duclos D. Contextual Note on the Ebola Bundibugyo Outbreak in Ituri (2026), SSHAP: \url{https://www.socialscienceinaction.org/resources/contextual-note-on-the-ebola-bundibugyo-outbreak-in-ituri-2026}
  \item Kheth'Impilo-DRC. Emergency Ebola outbreak response, Congo DR: project report, 1 Jul--31 Aug 2026 (B\'eni health zone): \url{https://cl.globalgiving.org/projects/emergency-ebola-outbreak-response-congo-dr/reports}
  \item de Padua B, Akhmetzhanov AR. Estimated transmissibility and case fatality rate of Bundibugyo virus, Uganda, 2007. Emerg Infect Dis. 2026;32(10):1762--1764. \url{https://wwwnc.cdc.gov/eid/article/32/10/26-1175}
  \item WHO. 10th Ebola virus disease outbreak in the DRC declared over, 25 Jun 2020: \url{https://www.who.int/news/item/25-06-2020-10th-ebola-outbreak-in-the-democratic-republic-of-the-congo-declared-over}
\end{itemize}

\section*{S3. Model and analysis code inventory}
\label{sec:si_code}

The analysis is fully reproducible from the deposited code. The core
modules (Python 3, \texttt{numpy}, \texttt{scipy}, \texttt{matplotlib}):

\begin{itemize}
  \item \texttt{ebola\_model\_v3.py} -- SEIHRD compartmental model with the path-mixed, care-seeking/community split and the explicit logistic detection-ramp observation layer; next-generation-matrix reproduction number and deterministic integrator.
  \item \texttt{v6\_extended.py} -- v6 calibration (V-C schedule), negative-binomial likelihood, Monte-Carlo propagation and one-year projection, ring-vaccination scenario module, the $k_{\mathrm{obs}}$ method-of-moments estimator (\texttt{mom\_k}) and the Monte-Carlo $k$ sampling set $\{12,\,18.44,\,39,\,\infty\}$.
  \item \texttt{\_v6\_paper\_quantities.py} -- the v6 paper quantities P1--P7 (calibration, identifiability, projection band, comparison table).
  \item \texttt{\_v6\_paper\_q123.py} -- supplementary quantities Q1--Q3 (completeness, multiplier bounds, scale factors).
  \item \texttt{\_v6\_mc\_boot.py} -- bootstrap / Monte-Carlo uncertainty propagation used for the projection intervals.
  \item \texttt{\_v6\_basins.py} -- maximized-likelihood basin and alternative-peak (B1--B4) diagnostics.
\end{itemize}

The full set of model, calibration, scenario and figure/table generation
scripts, together with the transcribed weekly series and data-processing
scripts, is deposited in the public version-controlled repository
archived with Zenodo (DOI \texttt{10.5281/zenodo.XXXXXXX}; replace with the
real deposited DOI before submission). No individual-level or unpublished
data were used; every plotted or tabulated number is produced directly
from the public series above.

\section*{S4. Reproducibility notes}
\label{sec:si_repro}

\begin{itemize}
  \item Environment: Python 3 (numpy, scipy, matplotlib); LaTeX with
        \texttt{amsmath}, \texttt{booktabs}, \texttt{xltabular}/\texttt{tabularx},
        \texttt{graphicx}, \texttt{hyperref}.
  \item Calibration: maximum-likelihood on the 23-interval likelihood
        (weekly UKHSA series + four daily register intervals), with the
        detection ceiling $d_{\max}=0.75$ fixed exogenously from the
        under-reporting anchor.
  \item To regenerate the manuscript numbers: run \texttt{v6\_extended.py}
        (calibration + propagation) then \texttt{\_v6\_paper\_quantities.py}
        and \texttt{\_v6\_paper\_q123.py}.
\end{itemize}

\end{document}
